\section{Falsification Criteria and Discriminant Tests}

A theory that cannot be falsified provides no scientific content. We state four explicit conditions under which the response-sector framework would be refuted.

\subsection{Criterion 1: Compression Failure}

If \textbf{(see PDF; target function/observable not recovered from DOCX extraction)} cannot be summarized by a universal low-DOF shape/amplitude tied to baryonic structure---if it requires an essentially free function per galaxy---then the ``materials reverse engineering'' story collapses into ad hoc curve-fitting and the framework is falsified as a parsimonious alternative.

\textbf{Status (DOCX-extracted; verify against PDF):} The single-parameter model achieves very strong improvement for 92\% of galaxies and passes BIC for 97\%, indicating compression holds across the SPARC sample.

\subsection{Criterion 2: Out-of-Sample Prediction Failure}

If internal predictors (baryonic structure, gas fraction, surface density) do not improve prediction of the edge amplitude in cross-validation beyond scale controls, the ``intrinsic boundary-layer'' causal story is not supported by the data.

\subsection{Criterion 3: Gravitational Slip / Lensing Inconsistency}

If the inferred extra potential needed for dynamics implies a (\textbf{see PDF; metric-potential pair not recovered}) pair that cannot match gravitational lensing constraints without adding new freedom, the metric-equivalence claim must be weakened or the framework revised. Specifically: if galaxy--galaxy lensing measurements require a different amplitude or profile than predicted by the response sector assuming \textbf{(see PDF; lensing/dynamics identification assumption)}, the model faces a decisive test (\S8).

\subsection{Criterion 4: Cosmological Growth Incompatibility}

If the response sector, extended to cosmological scales, produces a growth rate or perturbation structure that is incompatible with observed large-scale structure and CMB constraints, the framework cannot serve as a $\Lambda$CDM alternative (it may still be valid as a galaxy-domain model, but the ``go big'' program would be falsified).

\subsection{The Discriminant Triad}

If the framework is to serve as a serious $\Lambda$CDM-alternative candidate, three discriminant tests are paramount:
\begin{itemize}
	\item \textbf{Low-DOF compression across SPARC:} The response sector in the galaxy quasi-static limit is described by a universal family with \textbf{(see PDF; number of DOF per galaxy)} degrees of freedom per galaxy. Killer failure mode: needing an arbitrary function per galaxy.
	\item \textbf{Rotation vs.
	lensing consistency:} The response sector must give a definite relation between dynamics (\textbf{see PDF; rotation-curve observable}) and lensing (\textbf{see PDF; lensing observable}). If \textbf{(see PDF; no-slip condition)} in the relevant regime, state ``metric-equivalent in that regime.'' If not, predict where slip appears.
	\item \textbf{Cosmological growth and background viability:} Match background distances to first order and demonstrate that the response sector can cluster appropriately.
\end{itemize}
